\documentclass[12pt]{article}

\usepackage{geometry}
\geometry{a4paper, margin=2cm}

\usepackage{longtable}
\usepackage{caption}
\usepackage{array}
\usepackage{booktabs}
\usepackage[version=4]{mhchem}

% Column type with no hyphenation
\newcolumntype{L}[1]{>{\raggedright\arraybackslash}p{#1}}

\begin{document}

\section*{Scope of This Research — Table Test}

\subsection*{Table~\ref{tab:scope-full}: Scope boundary for the full research programme and its relation to future work.}

\setlength{\tabcolsep}{4pt}
\setlength{\LTleft}{0.025\textwidth}
\setlength{\LTright}{0.025\textwidth}
\begin{longtable}{|L{3.5cm}|L{2.5cm}|L{2.5cm}|L{7.0cm}|}
\hline
\textbf{Research Activity} & 
\textbf{In Scope} & 
\textbf{Out of Scope} & 
\textbf{How This Study Enables Future Work} \\
\hline
\endfirsthead

\hline
\textbf{Research Activity} & 
\textbf{In Scope} & 
\textbf{Out of Scope} & 
\textbf{How This Study Enables Future Work} \\
\hline
\endhead

\hline
\multicolumn{4}{|r|}{\textit{Continued on next page\ldots}} \\
\hline
\endfoot

\hline
\endlastfoot

% --- Core Reactor Characterisation ---
\multicolumn{4}{|L{14.5cm}|}{\textbf{Core Reactor Characterisation}} \\
\hline
Industry requirements matrix & 
Yes & --- & 
Defines power, range, and energy density targets that bound all downstream 
reactor design \\
\hline
Reaction parameter characterisation 
(T, P, particle size, fuel-to-water ratio) & 
Yes (Tier 1) & --- & 
Produces the experimental dataset for reactor geometry and operating point 
selection \\
\hline
By-product composition analysis 
(\ce{Al(OH)3}, AlOOH, \ce{Al2O3}) & 
Yes (Tier 1) & --- & 
Determines achievable effective energy density and recycling compatibility 
under real operating conditions \\
\hline
Seawater vs.\ distilled water reactant comparison & 
Yes (Tier 2) & --- & 
Establishes whether onboard seawater can be used directly, and specifies 
any pre-treatment or filtration requirements \\
\hline
Reactor start-up time and transient load characterisation & 
Yes (Tier 1) & --- & 
Provides the open-loop dynamic data required for control system design 
and battery buffer specification \\
\hline
Marine safety protocols (bunkering, dust control, emergency response) & 
Yes & --- & 
Provides the safety evidential basis for future certification submissions 
to RINA, LR, and DNV \\
\hline
Reactor geometry and material selection & 
Yes (lab-scale, indicative) & Full-scale design & 
Lab-scale findings provide the basis for scaled engineering design in 
future work \\
\hline
Preliminary cost estimation (reactor and aluminium fuel) & 
Yes (indicative, order-of-magnitude) & Final production costs & 
High-level estimates inform the economic feasibility case and identify 
where further cost refinement is needed \\
\hline

% --- Control Systems ---
\multicolumn{4}{|L{14.5cm}|}{\textbf{Control Systems}} \\
\hline
Closed-loop control system development & 
Yes (fixed-batch demonstrator) & Production-ready control hardware & 
Demonstrates thermal and kinetic controllability; quantifies response 
lags, rise times, and stability boundaries \\
\hline
Microcontroller hardware selection and implementation & 
Yes & Maritime ruggedisation & 
Provides a validated testbed for control algorithm development; hardware 
choices inform future procurement specifications \\
\hline
Power mode transition characterisation & 
Yes & Optimal controller design & 
Quantifies the reactor's dynamic response to step changes, enabling 
future engineers to size battery buffers and design control systems \\
\hline
Mode granularity optimisation & 
Yes & --- & 
Determines the optimal number and spacing of discrete power modes for 
marine propulsion, balancing controllability and manoeuvrability \\
\hline

% --- Engine and Battery Analysis ---
\multicolumn{4}{|L{14.5cm}|}{\textbf{Engine and Battery Analysis}} \\
\hline
Thermal integration pathway trade-off (direct vs. separation) & 
Yes & Detailed separation system design & 
Identifies the lowest-risk pathway for converting the steam-H$_2$ 
mixture into shaft power \\
\hline
Prime mover survey and selection & 
Yes (requirements definition) & Engine design or integration & 
Produces a decision matrix comparing mature, certifiable prime movers 
(steam turbine, ORC, diesel) \\
\hline
Battery buffer specification & 
Yes (requirements definition) & Battery system design & 
Defines minimum energy (kWh) and power (kW) requirements for a 
marine-certified battery to bridge reactor transients \\
\hline
Hybrid architecture recommendation & 
Yes (requirements definition) & Hybrid system integration & 
Recommends the simplest hybrid topology (series, parallel, or combined) 
that meets buffering requirements with lowest system risk \\
\hline
Efficiency target relaxation analysis & 
Yes & Efficiency optimisation & 
Establishes realistic, defensible efficiency targets (20--30\%) for a 
Generation-1 prototype \\
\hline

% --- Out of Scope ---
\multicolumn{4}{|L{14.5cm}|}{\textbf{Out of Scope (Future Work)}} \\
\hline
Full-scale reactor design and stress testing & 
--- & Yes & 
Lab-scale findings from this study provide the basis for scaling 
engineering design in future work \\
\hline
Continuous-flow reactor (real-time fuel injection \& solids handling) & 
--- & Yes & 
Fixed-batch control data demonstrates fundamental controllability; 
continuous-flow engineering is a separate mechanical design challenge \\
\hline
Engine and turbine detailed design and integration & 
--- & Yes & 
Interface specifications (steam temperature, pressure, flow rate) 
defined here enable future procurement and integration \\
\hline
Battery system detailed design and BMS development & 
--- & Yes & 
Capacity and power requirements defined here enable future procurement 
and integration \\
\hline
Classification society approval (full-scale) & 
--- & Yes & 
Safety protocols and experimental datasets developed here form the 
evidential basis for future certification submissions \\
\hline

\end{longtable}
\addtocounter{table}{-1}
\captionof{table}{Scope boundary for the full research programme and its relation to future work.}
\label{tab:scope-full}

\end{document}